\FloatBarrier
\section{Correlative microscopy}

Correlative microscopy is a generic term describing the combination of multiple microscopic and imaging methods.
In the field of biology, this term is well-established and has been used since at least 1956 \cite{bloch.1957}.
In this area, it usually represents the combination of light/fluorescence and electron microscopy \cite{bloch.1957,mironov.2013,loussertfonta.2015}.
Nonetheless, it can be applied to most microscopic and tomographic techniques, as long as the measurement is performed on exactly the same specimen with sufficient overlap \cite{loussertfonta.2015,daly.2017,mitchell.2019,mitchell.2021}.
Ideally, it is done within the same device, though this is not essential.

Advantages of the combination of multiple techniques are the combination of different information of one object/phase, such as:
\begin{itemize}
    \item trace element content (see \zcref{sec:corr_laicpms} and \cite{kleiner.2022b}),
    \item crystallographic information (see \zcref{sec:ebsd-hydrate} and \cite{rossler.2022}), and
    \item mechanical information (modulus, see \zcref{sec:EDX-NI} and \cite{chen.2010, li.2025, li.2026}).
\end{itemize}
Furthermore, information of a specimen across multiple scales can be collected \cite{daly.2017,mitchell.2019,mitchell.2021,kleiner.2022a}.
Finally, the resolution differences of different techniques can improve the segmentation results of lower resolution methods, enabling new insights, like to the properties of phases \cite{kleiner.2022b}.

%TODO: Du willst ja am Ende bzw. im Allgemeinen diskriminieren zwischen Phasen. Methode A hilft dir Phase A von B zu unterscheiden, Methode B bei der Unterscheidung zwischen B und C etc.

For example, the combination of \ac{SEM}-\ac{BSE}, \ac{EDX}, and \ac{EBSD} enables the identification of different phases by their chemical composition and the distinction between phases with varying crystal structures or orientations, even within a chemically identical area.

\subsection{Image alignment}

Good image alignment is the most important step before any further analysis.
While the automatic alignment of images with overlapping content from the same source (modality) is a mostly solved issue, multimodal image alignment is still a challenging problem \cite{frolov.2026}.
However, if the resulting images share enough features, they can still be aligned.
This can be done using
\begin{itemize}
    \item a fully manual alignment by rotating and transforming the images,
    \item a semi-manual method by identifying at least four matching point-pairs between both images and calculating the transformation \cite{hartley.2003},
    \item fully automatic methods from the field of computer vision/photogrammetry such as \ac{SIFT} \cite{lowe.2004} or machine learning-supported algorithms \cite{frolov.2025, frolov.2026}.
\end{itemize}

Usually, for few or single datasets the alignment can be done fully manually or by using tools like \textit{QGIS} \cite{QGIS_software}.
However, this can become tedious and time-consuming for more datasets, and in many cases, it is not very exact and not repeatable.

The main issues for the combination of multi-device (multi-modal) measurements are
\begin{itemize}
    \item distortions,
    \item different pixel aspect ratios,
    \item significant resolution differences,
    \item not enough reappearing matching features in one or both images (e.g. a phase which has a homogeneous chemical composition but differently oriented crystallites),
    \item and re-identifying the measurement positions in the first place.
\end{itemize}

% \color{orange}
% Aligning these images from multiple modalities can be helpful to find a fitting measurement position but also to evaluate the results in a later stage.
% The alignment can be done manually as discussed in the previous section by finding at least three or more matching point pairs or by manually transforming the images until the remaining error is acceptable.
% However, a more automatic method would be preferred.
% Multimodal image matching is a known but difficult problem in computer vision.
% Widley used algorithms like \ac{SIFT} usually expect images of the same modality and often fail if the modality changes.
% One possible solution is to bring all images in a similar modality.
% For example by first applying denoising and edge detection algorithms.

% \color{black}


% To simplify the realignment, markers (called fiducials) can be used to align different data sources to each other.
% However, this is not always possible or can be expensive.


